\chapter{Diseño}
\label{chap:deseño}

\lettrine{X}{xxx},\dots

 %%%%%%%%%%%%%%%%%%%%%%%%%%%%%%%%%%%%%%%%
 
\section{Introducción y objetivos}

\section{Resumen de patrones usados}

\section{Arquitectura general}

{ \color{red} Aspectos a considerar
    \begin{itemize}
      \item Arquitectura aplicación SPA
      \item División en subsistemas: En una aplicación SPA al menos backend y frontend pero podría haber más si se decide implementar utilidades e.g. como módulos independientes (e.g. como proyectos maven / jars independientes)
    \end{itemize}
}
%%%%%%%%%%%%%%%%%%%%%%%%%%%%%%%%%%%%%%%%
% Subsistema Backend %
%%%%%%%%%%%%%%%%%%%%%%%%%%%%%%%%%%%%%%%%
\section{Subsistema Backend}

\subsection{Objetivos}

\subsection{Arquitectura}

{ \color{red} Aspectos a considerar
    \begin{itemize}
      \item Diagrama de paquetes general del subsistema (figura y descripción de la misma)
    \end{itemize}
}

\subsection{Modelo del dominio}

\subsubsection{Diagrama de entidades}

{ \color{red} Aspectos a considerar
    \begin{itemize}
      \item Diagrama de clases persistentes de la aplicación
      \item Para cada indicar, tanto de su semántica como de sus atributos.
      \item Justificar decisiones de diseño (¿desnormalizaciones?, métodos de lógica - Domain Model Pattern - ?)
    \end{itemize}
}
\subsubsection{Modelo de datos}

{ \color{red} Aspectos a considerar
    \begin{itemize}
      \item Diagrama entidad-relación
      \item Diccionario de Datos (en su ausencia - por temas de espacio bastante razonable - compensable con defición clara en diagrama de entidades).
    \end{itemize}
}

\subsubsection{Consideraciones adicionales de modelado de datos}

{ \color{red} Aspectos a considerar
    \begin{itemize}
      \item Claves primarias autogeneradas
      \item Anotaciones para relaciones entre entidades
      \item Optimistic Locking y @Version en Hibernate
      \item Relaciones LAZY/EAGER
      \item Optimizaciones hibernate en navegación entre entidades
      \item Optimizaciones hibernate - cache primer nivel
    \end{itemize}
}

\subsection{Capa de acceso a datos}

{ \color{red} Aspectos a considerar
    \begin{itemize}
      \item Arquitectura de un DAO
      \item Aspectos más relevantes de los DAOs implementados
    \end{itemize}
}

\subsection{Capa de lógica de negocio}

{ \color{red} Aspectos a considerar
    \begin{itemize}
      \item Arquitectura de servicios del modelo: Interfaces + clases asociadas
      \item Incluir descripción de métodos (en algunos casos podrían comentarse de forma agregada algunas operaciones por ejemplo gestión de usuarios (añadir/eliminar/modificar) en lugar de comentar de uno en uno.
      \item Incluir descripción de algoritmos en los casos que sea necesario. Podría incluirse un diagrama de secuencia para complementar algún algoritmo.
    \end{itemize}
}

\subsection{Capa servicios}

{ \color{red} Aspectos a considerar
    \begin{itemize}
      \item Diseño controladores
      \item API REST (intentar que sea REST-full cuando haya un mapping directo ... overloading del POST en otro caso
      \item Gestión de errores y aspectos de internacionalización de mensajes de error
      \item Autenticación y aspectos de control de acceso a los diferentes endpoints (SecurityConfig). Soporte de redirección tras autenticación, si implementado.
    \end{itemize}
}

\subsection{Otros?}

{ \color{red} Aspectos a resaltar
    \begin{itemize}
      \item Integración con otros sistemas o librerías a nivel de backend, e.g. pago con PayPal. 
      \item Tb se podrían destacar aspectos como chats con websockets o algoritmos. Importante dejar claro todo lo que se comente en "objetivos" y "aportaciones" principales del proyecto.
    \end{itemize}
}

%%%%%%%%%%%%%%%%%%%%%%%%%%%%%%%%%%%%%%%%
% Subsistema Frontend %
%%%%%%%%%%%%%%%%%%%%%%%%%%%%%%%%%%%%%%%%
\section{Subsistema Frontend}

\subsection{Objetivos}

\subsection{Arquitectura}

{ \color{red} Aspectos a considerar
    \begin{itemize}
      \item Arquitectura de capas: al menos capa de acceso al servicio y capa de componentes de interfaz
    \end{itemize}
}

\subsection{Capa de acceso al servicio}

\subsection{Capa de componentes de interfaz}

\subsection{Otros?}

{ \color{red} Aspectos a considerar
    \begin{itemize}
      \item Internacionalización de mensajes de texto 
      \item Integración con otros sistemas / librerías
    \end{itemize}
}
